\section{Proofs for the exact number of modes}\label{app:modes}

This appendix proves parts (iii) to (v) of Theorem~\ref{thm:mode-count},
Lemma~\ref{lem:diff-bigjump} and Corollary~\ref{cor:sequences}, and checks the junctions of the
five zones. Parts (i) and (ii) are Theorem~\ref{thm:ends}.

\emph{Notation.} For a sequence $X$ on $\Z$ we write $\Delta X(j)=X(j+1)-X(j)$ and
$\nabla X(n)=X(n-1)-X(n)$, and $\Delta^2$, $\nabla^2$ for the iterates. For a family $X_m$
indexed by $m$ we write $\delta X_m=X_{m+1}-X_m$ and $\delta^2X_m=\delta X_{m+1}-\delta X_m$.
For real $x\ge1$ we read $f_\sgn(x)$ as $f_\sgn(\lfloor x\rfloor)$. By Lemma~\ref{lem:rv}(ii)
and $mT_\sgn(a_\sgn(m))\to1$, for every $\varepsilon'>0$ there are $c,C>0$ with
\begin{equation}\label{eq:potter-m}
 c\,y^{-\alpha_\sgn-\varepsilon'}\le m\,T_\sgn\bigl(y\,a_\sgn(m)\bigr)\le C\,y^{-\alpha_\sgn+\varepsilon'}
 \qquad(m\ge1,\ y\ge1).
\end{equation}
For bounded $m$ this follows from the bounds for $T_\sgn$ alone, since $a_\sgn(m)$ is then
bounded above and below. Below $y=1$ the exponents exchange. With $x_0$ the threshold of
Potter's bounds in Lemma~\ref{lem:rv}(ii), for every $\varepsilon'>0$ there are $c,C>0$ with
\begin{equation}\label{eq:potter-m-small}
 c\,y^{-\alpha_\sgn+\varepsilon'}\le m\,T_\sgn\bigl(y\,a_\sgn(m)\bigr)\le C\,y^{-\alpha_\sgn-\varepsilon'}
 \qquad(m\ge1,\ 0<y\le1,\ y\,a_\sgn(m)\ge x_0),
\end{equation}
since then both $a_\sgn(m)$ and $ya_\sgn(m)$ are at least $x_0$ and $mT_\sgn(a_\sgn(m))$ is
bounded above and below.

\subsection{Smoothness and the renewal sequence}\label{app:smooth}

\begin{lemma}[consequences of (LD2)]\label{lem:ld2}
Let $f$ be a positive sequence that is regularly varying with index $-\alpha-1$, $\alpha>0$, and
satisfies $|f(n+1)-2f(n)+f(n-1)|\le Cf(n)/n^2$. Then $|f(n+1)-f(n)|\le C'f(n)/n$ for large $n$,
and $f(n)-f(n+1)=(\alpha+1+o(1))f(n)/n$.
\end{lemma}

So under (LR), (LD2) implies (LD), and $\nabla f_\sgn(y)=(\alpha_\sgn+1+o(1))f_\sgn(y)/y$ as
$y\to\infty$.

\begin{proof}
Since $f(n)\to0$, $\Delta f(n)\to0$, and $\Delta f(n)=-\sum_{i\ge n}\Delta^2f(i)$. Here
$|\Delta^2f(i)|\le Cf(i+1)/(i+1)^2\le C'f(i)/i^2$, and $\sum_{i\ge n}f(i)/i^2\le C''f(n)/n$ by
Karamata's theorem, since $f(i)/i^2$ has index $-\alpha-3$. This is the first claim. For the
second put $r(n)=n\Delta f(n)/f(n)$, which is bounded. Fix $\lambda>1$ and let
$k=\lfloor(\lambda-1)n\rfloor$. For $0\le i<k$ the bound on $\Delta^2f$ gives
$|\Delta f(n+i)-\Delta f(n)|\le C_\lambda if(n)/n^2$, with $C_\lambda$ bounded for
$\lambda\le2$. Summing over $i<k$,
\[
 \frac{f(n+k)}{f(n)}-1=\frac kn\,r(n)+O\bigl((\lambda-1)^2\bigr).
\]
The left side tends to $\lambda^{-\alpha-1}-1$ by Lemma~\ref{lem:rv}(i), and $k/n\to\lambda-1$.
Hence $\limsup_n|r(n)-(\lambda^{-\alpha-1}-1)/(\lambda-1)|\le C(\lambda-1)$. The quotient tends to
$-(\alpha+1)$ as $\lambda\downarrow1$, so $r(n)\to-(\alpha+1)$.
\end{proof}

\begin{lemma}[second-difference renewal theorem]\label{lem:second-renewal}
Assume (RV), (SPAN) and (LR) for the minus sign, and write $u=u_\minus$, $v=v_\minus$,
$t_n=T_\minus(n+1)$ and $d_n=u(n)-u(n-1)$ as in Lemma~\ref{lem:diff-renewal}. Then
$n(d_n-d_{n-1})/t_n\to\alpha_\minus/\mu_\minus^2$. Equivalently
\[
 \Delta^2u(j)=u(j+2)-2u(j+1)+u(j)=\bigl(1+o(1)\bigr)\frac{f_\minus(j)}{\mu_\minus^2}.
\]
Moreover $|\Delta^2u(j)|\le CT_\minus(j)/(j+1)$ for $j\ge0$, $\Delta^2v(j)=O(T_\minus(j))$, and the
sequence $\tilde v$ equal to $v$ on $\Z_{\ge0}$ and to $0$ elsewhere has $\Delta^2\tilde v\in\ell^1$.
\end{lemma}

\begin{proof}
We repeat the proof of Lemma~\ref{lem:diff-renewal} with one more seminorm. For sequences $a$
on $\Z_{\ge0}$ put $\nabla a_n=a_n-a_{n-1}$ with $a_{-1}=0$, and
$\|a\|_{t'}=\sup_{n\ge1}n|\nabla a_n|/t_n$. For $n\ge2$ and $M=\lfloor n/2\rfloor$ a direct
computation gives
\[
 \nabla(a*b)_n=\sum_{k\le M}a_k\nabla b_{n-k}+\sum_{l\le n-M-2}b_l\nabla a_{n-l}+a_{M+1}b_{n-M-1}.
\]
In the two sums the index of the differenced factor is at least $n/2$, and the last term is at
most $\|a\|_t\|b\|_tc_t'^2t_n^2$, where $c'_t=\sup_nt_{\lfloor n/2\rfloor-1}/t_n$. Since
$nt_n\to0$,
\begin{equation}\label{eq:weighted2}
 \|a*b\|_{t'}\le C\bigl(\|a\|_1\|b\|_{t'}+\|a\|_{t'}\|b\|_1+\|a\|_t\|b\|_t\bigr),
\end{equation}
where the constant of the first two terms is $2c_t$. If moreover $n\nabla a_n/t_n\to\chi'_a$ and
$n\nabla b_n/t_n\to\chi'_b$, then dominated convergence in the two sums, as
for~\eqref{eq:product-rule}, gives $n\nabla(a*b)_n/t_n\to\hat a(1)\chi'_b+\hat b(1)\chi'_a$.

By (LR), $n\nabla t_n/t_n=-nf_\minus(n)/T_\minus(n+1)\to-\alpha_\minus$, and the same holds for
the truncation $t''$ of the proof of Lemma~\ref{lem:diff-renewal}. Take $M$ there so large that
also $2c_t\|g\|_1\le\frac14$. This is possible because
$\|e\|_1\le2\|d\|_1$ for every $M$. The coefficients $e$ decay geometrically, so
$n\nabla e_n/t_n\to0$, and $g=e*t''$ has $n\nabla g_n/t_n\to-\alpha_\minus\hat e(1)$.
By~\eqref{eq:weighted2} and the bound $\|g^{*k}\|_t\le Ck4^{-k}$ of that proof,
$\|g^{*k}\|_{t'}\le Ck^24^{-k}$. By induction
$n\nabla(g^{*k})_n/t_n\to k\hat g(1)^{k-1}(-\alpha_\minus\hat e(1))$, and dominated convergence
over $k$ gives, for $h=\sum_{k\ge1}(-g)^{*k}$, the limit $\alpha_\minus\hat e(1)/(1+\hat g(1))^2$
of $n\nabla h_n/t_n$. Since $d=e+e*h$, the limit of $n\nabla d_n/t_n$ is
$\hat e(1)\alpha_\minus\hat e(1)/(1+\hat g(1))^2=\alpha_\minus/\mu_\minus^2$, with
$\hat e(1)/(1+\hat g(1))=1/\mu_\minus$ as in that proof. By (LR) this is the stated form of
$\Delta^2u(j)=\nabla d_{j+2}$, and $|\nabla d_n|\le Ct_n/n$ gives the bound on $\Delta^2u$.

The generating function of $d*d$ is $(1-z)^2\hat u^2$, so $d*d=\nabla^2(u*u)$, and
$(d*d)_n/t_n\to2\hat d(1)(-1/\mu_\minus^2)$ by~\eqref{eq:product-rule}. By~\eqref{eq:conv-v},
$u*u=v+u$ on $\Z_{\ge0}$. Hence $\nabla^2v_n=(d*d)_n-\nabla^2u_n=O(t_n)$, that is
$\Delta^2v(j)=O(T_\minus(j))$. Since $T_\minus$ is summable, $\Delta^2\tilde v\in\ell^1$.
\end{proof}

\subsection{One antimode}\label{app:antimode}

Write $f=f_\plus$, $T=T_\plus$, $\mu=\mu_\plus$, $V_m=u^\plus_m$ with $V_0=\delta_0$ and
$q_m=q^\plus_m$, as in Lemma~\ref{lem:plus-diff}.

\begin{lemma}[plus runs at second order]\label{lem:plus-second}
Assume (LR) and (LD2) for the plus sign, and fix $\eta\in(0,1)$. As $n\to\infty$, uniformly in
$0\le m\le n^{1-\eta}$,
\begin{enumerate}[(i)]
\item $|\nabla^2V_m(n)|\le Cmf(n)/n^2$,
\item $|\nabla\delta V_m(n)|\le Cf(n)/n$,
\item $|\delta^2V_m(n)|\le Cf(n)/n$,
\item $|q_{m+3}(n)-2q_{m+2}(n)+q_{m+1}(n)|\le Cf(n)/n$.
\end{enumerate}
\end{lemma}

\begin{proof}
We use Lemma~\ref{lem:plus-diff} and its proof freely, with $\eta/2$ in place of $\eta$ at
arguments $n'\ge n/2$. In particular $V_m(x)\le Cmf(n)$ for $m\ge1$ and $x\ge n/2$, and
$mnf(n)\le Cn^{1-\eta}T(n)\to0$. By Lemma~\ref{lem:ld2}, (LD) holds for the plus sign.

(i) For $m=0$ there is nothing to prove. Let $m\ge1$, fix $K_0\ge4$ with
$(K_0/4-1)(\alpha_\plus-1)/2>\alpha_\plus+4$, and put $w=\lceil n/K_0\rceil$. As in the proof of
Lemma~\ref{lem:plus-diff}(ii), Lemma~\ref{lem:truncated} bounds the probability that
$S^\plus_m=n'$ with no summand of size at least $w$ by $Cn^{-\alpha_\plus-4}=o(f(n)/n^2)$, for each
$n'\in\{n-2,n-1,n\}$. Otherwise let $I$ be the first index with $X^\plus_I\ge w$ and let $h_i$ be
as in that proof. The second difference of
$\Prob(S^\plus_m=n',I=i)=\sum_zh_i(z)f(n'-z)\ones_{\{n'-z\ge w\}}$ at $n$ is
\[
 \sum_{z\le n-2-w}h_i(z)\nabla^2f(n-z)+h_i(n-1-w)\bigl(f(w+1)-f(w)\bigr)+f(w)\bigl(h_i(n-w)-h_i(n-1-w)\bigr).
\]
In the first sum $n-z\ge n/K_0$, so (LD2) and Lemma~\ref{lem:rv}(ii) bound it by $Cf(n)/n^2$. In
the second term $h_i(n-1-w)\le V_{m-1}(n-1-w)\le Cmf(n)$ and $|f(w+1)-f(w)|\le Cf(n)/n$ by (LD).
For the third term, $h_i$ is the law of a sum of $m-1$ independent variables, $i-1$ of them
with the defective law $f\ones_{\{\cdot<w\}}$. We apply the proof of
Lemma~\ref{lem:plus-diff}(ii) to this sum at the level $n-w\ge n/2$ with the threshold $w$
itself. A truncated variable is then never a giant summand, the bound with no giant summand is
the same as before, and the giant summands are among the untruncated variables. This gives
$|h_i(n-w)-h_i(n-1-w)|\le Cmf(n)/n+Cm^2f(n)^2\le Cmf(n)/n$. Summing over $i\le m$, the second
difference is at most $Cmf(n)/n^2+Cm^2f(n)^2/n\le Cmf(n)/n^2$.

(ii) For $m=0$, $\nabla\delta V_0(n)=\nabla f(n)$ for $n\ge2$. Let $m\ge1$. Since
$\delta V_m=f*V_m-V_m$,
\[
 \nabla\delta V_m(n)=\sum_{x\le n/2}f(x)\bigl[\nabla V_m(n-x)-\nabla V_m(n)\bigr]
 -\nabla V_m(n)\sum_{x>n/2}f(x)+\sum_{y<n/2}f(n-y)\nabla V_m(y).
\]
By (i) the first bracket is at most $Cxmf(n)/n^2$, so the first sum is at most
$C\mu mf(n)/n^2\le Cf(n)/n$. The second term is at most $Cmf(n)T(n/2)/n$. In the last sum we sum
by parts in $y$. This moves the difference onto $f(n-\cdot)$, where (LD) bounds it by $Cf(n)/n$,
and leaves boundary terms of the form $f(n-y)V_m(y)$ with $y\ge n/2-1$, which are at most
$Cmf(n)^2$. Since $\sum_yV_m(y)\le1$, the last sum is at most $Cf(n)/n$.

(iii) Since $\delta^2V_m=f*\delta V_m-\delta V_m$,
\[
 \delta^2V_m(n)=\sum_{x\le n/2}f(x)\bigl[\delta V_m(n-x)-\delta V_m(n)\bigr]
 -\delta V_m(n)\sum_{x>n/2}f(x)+\sum_{y<n/2}f(n-y)\delta V_m(y).
\]
By (ii) the first sum is at most $C\mu f(n)/n$. The second term is $O(f(n)T(n))$ by
Lemma~\ref{lem:plus-diff}(iii). Put $\gamma(y)=(f(n-y)-f(n))\ones_{\{0\le y<n/2\}}$. The last sum
equals $f(n)\sum_{y<n/2}\delta V_m(y)+\E[\gamma(S^\plus_m+X)-\gamma(S^\plus_m)]$, with $X$ an
independent copy of $X_\plus$. Here $\sum_{y<n/2}\delta V_m(y)=-q_{m+1}(\lceil n/2\rceil)=O(T(n))$
by Lemma~\ref{lem:plus-diff}(i). In the expectation, if both arguments are below $n/2$ the
difference is at most $CXf(n)/n$ by (LD). If only the first one is, it is at most $Cf(n)$, and
this has probability $q_{m+1}(\lceil n/2\rceil)\le CT(n)$. Since $nT(n)\to0$, the total is at
most $Cf(n)/n$.

(iv) Write $\delta q_{m+1}=q_{m+2}-q_{m+1}$. As in the proof of Lemma~\ref{lem:plus-diff}(iv),
$\delta q_{m+1}(n)=\sum_xf(x)[q_{m+1}(n-x)-q_{m+1}(n)]$ and $q_{m+1}(l-1)-q_{m+1}(l)=\delta V_m(l-1)$.
Hence $\nabla\delta q_{m+1}(l)=\delta^2V_m(l-1)$, and the quantity in (iv) is
$\sum_xf(x)[\delta q_{m+1}(n-x)-\delta q_{m+1}(n)]$. We split it as in (iii). The part
$x\le n/2$ is at most $C\mu f(n)/n$ by (iii), and the part with $\sum_{x>n/2}f(x)$ is $O(f(n)T(n))$
by Lemma~\ref{lem:plus-diff}(iv). The rest is
$f(n)\sum_{y<n/2}\delta q_{m+1}(y)+\sum_y\gamma(y)\delta q_{m+1}(y)$. Since
$\delta q_{m+1}=f*q_{m+1}-q_{m+1}$,
\[
 \sum_y\gamma(y)\delta q_{m+1}(y)=\sum_xf(x)\sum_{y'}q_{m+1}(y')\bigl[\gamma(y'+x)-\gamma(y')\bigr].
\]
The terms with $y'+x<n/2$ add up to at most $C\mu\sum_{y'}q_{m+1}(y')f(n)/n=C\mu^2f(n)/n$. The terms
with $y'<n/2\le y'+x$ are at most $Cf(n)\sum_{y'<n/2}q_{m+1}(y')T(\lceil n/2\rceil-y')$, and this
sum is at most $C\mu T(n)$. Indeed, for $y'\le n/4$ the factor $T$ is at most $T(n/4)$ and the
weights add up to $\mu$, and for $y'>n/4$ the weight $q_{m+1}(y')$ is at most $CT(n)$ by
Lemma~\ref{lem:plus-diff}(i) while the factors $T$ add up to at most $\mu$. Finally the full sum
of $\delta q_{m+1}$ vanishes, and for every $L'$
\[
 \Bigl|\sum_{y\ge L'}\delta q_{m+1}(y)\Bigr|=\sum_xf(x)\sum_{L'-x\le y<L'}q_{m+1}(y)\le(C+1)\mu\,T(L'/2),
\]
by the same split at $x=L'/2$. With $L'=\lceil n/2\rceil$ this part is $f(n)O(T(n))$.
\end{proof}

\begin{theorem}[second differences in the layer]\label{thm:layer-second}
Assume (RV), (SPAN) and (LR) for the minus sign and (LR) and (LD2) for the plus sign, and fix
$\eta\in(0,1)$. Let the start be fresh with either sign, or a mixture of the two fresh starts.
Then, uniformly in $1\le j\le N^{1-\eta}$,
\[
 D_N(j+1)-D_N(j)=T_\plus(N)\,\Delta^2u_\minus(j)+O\Bigl(f_\plus(N)\Bigl(T_\minus(j)+\frac1N\Bigr)\Bigr).
\]
\end{theorem}

\begin{proof}
By linearity we may treat the two fresh starts. We keep the notation of the proof of
Theorem~\ref{thm:first-diff}, with $n=N-j$, $a_m=u^\minus_m$, and $Q_m$, $b_m$ as there, and we
use its six facts. Differencing $D_N=\mathrm I+\mathrm{II}+\mathrm{III}+\mathrm{IV}$ once more in
$j$, with $n$ replaced by $n-1$, gives the exact identity
\[
 D_N(j+1)-D_N(j)=A_1+2A_2+A_3+2A_4+A_5+A_6
\]
for a plus start, where
\begin{align*}
 A_1&=\sum_m\Delta^2a_m(j)\,Q_m(n-2),& A_2&=\sum_m\Delta a_m(j)\,\delta V_m(n-2),\\
 A_3&=\sum_ma_m(j)\,\nabla\delta V_m(n-1),& A_4&=\sum_m\nabla V_m(n-1)\,\Delta b_m(j),\\
 A_5&=\sum_m\nabla^2V_m(n)\,b_m(j),& A_6&=\sum_mV_m(n-2)\,\Delta^2b_m(j).
\end{align*}
Here we used $\nabla Q_m(l)=\delta V_m(l-1)$. Only $m\le j+2$ contributes, so
Lemma~\ref{lem:plus-second} applies at the arguments $n-2\ge N/2$ with $\eta/2$ in place of
$\eta$, and $f(n')\le Cf(N)$ there.

By Lemma~\ref{lem:plus-second}(i) and fact (6), $|A_5|\le Cf_\plus(N)j/N^2$. By (ii) and fact
(1), $|A_3|\le Cf_\plus(N)/N$. Writing $\nabla V_m=\nabla V_1+\sum_{i<m}\nabla\delta V_i$ and using
$\sum_m\Delta b_m(j)=0$ and fact (5), $A_4=\sum_{i\ge1}\nabla\delta V_i(n-1)u^\minus_i(j)$, which
is at most $Cf_\plus(N)/N$. In the same way $A_6=\sum_{i\ge1}\delta V_i(n-2)\Delta a_i(j)$, which
has the form of $A_2$. Summation by parts with $\delta V_m=\delta V_1+\sum_{i<m}\delta^2V_i$ gives
$A_2=\delta V_1(n-2)\Delta u(j)+\sum_{i\ge1}\delta^2V_i(n-2)\Delta\bar A_i(j)$. By
Lemma~\ref{lem:plus-diff}(iii), Lemma~\ref{lem:diff-renewal}, Lemma~\ref{lem:plus-second}(iii)
and fact (2), $A_2=O(f_\plus(N)T_\minus(j))+O(f_\plus(N)/N)$.

For $A_1$ write
$Q_m=Q_1+(m-1)\delta Q_1+\sum_{i=1}^{m-2}(m-1-i)\delta^2Q_i$ and sum by parts twice. With
$\sum_ma_m=u$, $\sum_m(m-1)a_m=v-u$ and
\[
 W_i(j)=\sum_{m\ge i+2}(m-1-i)\,a_m(j)=\sum_xu^\minus_{i+1}(x)\,\tilde v(j-x),
\]
this gives
\[
 A_1=Q_1(n-2)\Delta^2u(j)+\delta Q_1(n-2)\Delta^2(v-u)(j)+\sum_{i\ge1}\delta^2Q_i(n-2)\Delta^2W_i(j).
\]
Here $Q_1(n-2)=q_2(N-j-2)=T_\plus(N)+O((j+1)f_\plus(N))$ by Lemma~\ref{lem:plus-diff}(iv), and
$(j+1)|\Delta^2u(j)|\le CT_\minus(j)$ by Lemma~\ref{lem:second-renewal}. So the first term is
$T_\plus(N)\Delta^2u(j)+O(f_\plus(N)T_\minus(j))$. In the second, $\delta Q_1=O(f_\plus(N))$ by
Lemma~\ref{lem:plus-diff}(iv) and $\Delta^2(v-u)(j)=O(T_\minus(j))$ by
Lemma~\ref{lem:second-renewal}. In the third, $|\delta^2Q_i(n-2)|\le Cf_\plus(N)/N$ by
Lemma~\ref{lem:plus-second}(iv), and
$\sum_i|\Delta^2W_i(j)|\le\sum_xu(x)|\Delta^2\tilde v(j-x)|\le\|\Delta^2\tilde v\|_1$. Adding up
gives the theorem for a plus start.

For a minus start $Q_m=q^\plus_m$ and $b_m=q^\minus_{m+1}$. The index of every $\delta V$ and
$\delta^2Q$ above drops by one, which Lemmas~\ref{lem:plus-diff} and~\ref{lem:plus-second} allow,
since they include $m=0$. Now $Q_1=T_\plus$, $\sum_m\Delta b_m(j)=f_\minus(j)$ and
$\sum_m\Delta^2b_m(j)=\Delta f_\minus(j)$. This adds the terms $\nabla V_1(n-1)f_\minus(j)=O(f_\plus(N)/N)$
to $A_4$ and $V_1(n-2)\Delta f_\minus(j)=O(f_\plus(N)T_\minus(j))$ to $A_6$. The rest is unchanged.
\end{proof}

\begin{corollary}[convexity near the plus end]\label{cor:convex}
Assume the hypotheses of Theorem~\ref{thm:layer-second} with $\eta<1-1/\alpha_\minus$. Let $J_0$
be such that $f_\minus>0$ and $\Delta^2u_\minus>0$ on $[J_0,\infty)$, which exists by (LR) and
Lemma~\ref{lem:second-renewal}. For every $K$ and all large $N$, $D_N$ is strictly increasing on
$J_0\le j\le K\Lambda_N$.
\end{corollary}

\begin{proof}
By Lemma~\ref{lem:second-renewal} there is $c>0$ with $\Delta^2u_\minus(j)\ge cf_\minus(j)$ for
$j\ge J_0$. Since $f_\plus(N)\le CT_\plus(N)/N$, it suffices that $Nf_\minus(j)/T_\minus(j)\to\infty$
and $N^2f_\minus(j)\to\infty$ uniformly on $[J_0,K\Lambda_N]$, which lies in $[1,N^{1-\eta}]$ for
large $N$. The first ratio is at least $cN/j\ge cN/(K\Lambda_N)$, by (LR). For the second,
Lemma~\ref{lem:rv}(ii) gives $f_\minus(j)\ge cf_\minus(K\Lambda_N)$ on the range, and
$N^2f_\minus(K\Lambda_N)\ge cN\cdot NT_\minus(K\Lambda_N)/(K\Lambda_N)\ge cK^{-\alpha_\minus-1}N/\Lambda_N$
for large $N$. Both tend to infinity since $\Lambda_N=o(N)$.
\end{proof}

\begin{proof}[Proof of Theorem~\ref{thm:mode-count}(iii)]
For large $N$ the window lies in $[J_0,(x_*+1)\Lambda_N]$. Apply Corollary~\ref{cor:convex} with
$K=x_*+1$. By Theorem~\ref{thm:ends}(d) $D_N$ changes sign in the window, and a strictly
increasing sequence changes sign at most once and vanishes at most once. (EQ) was not used.
\end{proof}

\subsection{Fourier bounds and truncated sums}\label{app:fourier}

In this subsection we fix a sign and drop it, and we assume (RV) and (SPAN). So $f$, $T$, $\mu$,
$\alpha$, $a$, $S_m$ and $u_m=f^{*m}$ refer to that sign, and $\varphi(t)=\E e^{itX}$. We use
three elementary bounds. First, for all $m\ge1$ and $y>0$,
\begin{equation}\label{eq:right-tail}
 \Prob(S_m-m\mu\ge y)\le C\,mT(y).
\end{equation}
If $mT(y/8)\ge1$ this holds since $T(y/8)\le C_8T(y)$. Otherwise either some summand is at least
$y/8$, which has probability at most $mT(y/8)$, or none is, and Lemma~\ref{lem:truncated} with
$w=y/8$ bounds this case by $e^{C_{\mathrm K}}(mT(y/8))^4$. Second, for $m\ge1$ and $y\ge0$,
\begin{equation}\label{eq:left-local}
 u_m(m\mu-y)\le\frac C{a(m)}\,e^{-y/(2a(m))}.
\end{equation}
For $m=1$ the left side vanishes unless $y<\mu$. For $m\ge2$ write $u_m=u_{m_1}*u_{m_2}$ with
$m_1=\lfloor m/2\rfloor$. One of the two blocks lies at least $y/2$ below its mean, so the left
side is at most $\sup u_{m_2}\Prob(S_{m_1}\le m_1\mu-y/2)+\sup u_{m_1}\Prob(S_{m_2}\le m_2\mu-y/2)$,
and Lemma~\ref{lem:stable}(ii) and~(iv) conclude, since $a(m_i)\le a(m)\le Ca(m_i)$. Third,
$\E|S_m-m\mu|\le Ca(m)$, since by~\eqref{eq:right-tail} and~\eqref{eq:karamata}
$\E(S_m-m\mu)_+\le a(m)+CmC_{\mathrm K}a(m)T(a(m))$, and the negative part has the same mean.

We also need $1-|\varphi(t)|\ge cT(1/|t|)$ for $|t|\le\pi$. For small $|t|$ this follows from
$1-|\varphi(t)|^2=\E[1-\cos t(X-X')]\ge\frac14t^2\E[(X-X')^2;|X-X'|\le1/|t|]$, with $X'$ an
independent copy, and from $\Prob(X-X'\in(y/2,y])\ge f(n_0)\Prob(X-n_0\in(y/2,y])\ge cT(y)$ for
any $n_0$ with $f(n_0)>0$. For the other $t$ it follows from (SPAN), since $|\varphi|<1$ on
$0<|t|\le\pi$.

\begin{lemma}[truncated sums]\label{lem:trunc-local}
For $w\ge1$ let $u^{(w)}_m(x)=\Prob(S_m=x,\ \max_{i\le m}X_i<w)$.
\begin{enumerate}[(a)]
\item For all $m\ge1$, $w\ge1$ and integers $x$ with $x-m\mu\ge y\ge2w$,
$u^{(w)}_m(x)\le Ca(m)^{-1}(mT(w))^{y/(4w)}$ whenever $mT(w)\le1$.
\item If $mT(w)\le1$, then $\sup_x|\nabla u^{(w)}_m(x)|\le Ca(m)^{-2}$.
\end{enumerate}
\end{lemma}

\begin{proof}
(a) For $m=1$ the left side vanishes, since $x\ge y+\mu>w$. For $m\ge2$ split $m=m_1+m_2$ as
above. One of the two blocks exceeds its mean by at least $y/2$, so
$u^{(w)}_m(x)\le\sum_{i=1,2}\sup u_{m_{3-i}}\cdot\Prob(S_{m_i}-m_i\mu\ge y/2,\ \max<w)$, and
Lemma~\ref{lem:stable}(ii) and Lemma~\ref{lem:truncated} give the bound.

(b) The generating function of $u^{(w)}_m$ is $\varphi_w^m$ with
$\varphi_w(t)=\sum_{n<w}f(n)e^{itn}$, and $|\varphi_w|\le|\varphi|+T(w)$. Hence
$|\varphi_w(t)|^m\le e^{mT(w)}e^{-cmT(1/|t|)}\le e\cdot e^{-cmT(1/|t|)}$ and
\[
 |\nabla u^{(w)}_m(x)|\le\frac e{2\pi}\int_{-\pi}^\pi|t|\,e^{-cmT(1/|t|)}\,dt .
\]
For $m\ge m_0$ substitute $t=s/a(m)$. By~\eqref{eq:potter-m-small} with $y=1/s$,
$mT(a(m)/s)\ge cs^{\alpha-\varepsilon'}$ for $1\le s\le a(m)/x_0$, and for $|t|\ge1/x_0$ the integrand is at most $e^{-cmT(x_0)}$. So the
integral is at most $Ca(m)^{-2}$. For $m<m_0$ the integral is bounded and $a(m)$ is bounded.
\end{proof}

\begin{lemma}[local limits for differences]\label{lem:fourier-diff}
Let $\mathsf k$ be a kernel of the fixed sign, that is $u_m$, $q_m$ or $q_{m+1}$, with weight
$c_{\mathsf k}$ as in Appendix~\ref{sec:bulk-kernels}, and let $r\in\{1,2\}$. Then
\[
 \sup_{j\in\Z}\Bigl|a(m)^{1+r}\Delta^r\mathsf k_m(j)-c_{\mathsf k}\,g^{(r)}\Bigl(\frac{j-m\mu}{a(m)}\Bigr)\Bigr|\to0
 \quad(m\to\infty),
\]
and $\sup_{m\ge1}\sup_ja(m)^{1+r}|\Delta^r\mathsf k_m(j)|<\infty$.
\end{lemma}

\begin{proof}
This is Gnedenko's proof~\cite[Theorem~4.2.1]{ibragimovlinnik1971} with one more factor. We have
$\Delta^ru_m(j)=(2\pi)^{-1}\int_{-\pi}^\pi\varphi(t)^m(e^{-it}-1)^re^{-itj}\,dt$. On
$\varepsilon\le|t|\le\pi$, (SPAN) gives $|\varphi|\le e^{-c_\varepsilon}$, so this range
contributes $O(e^{-c_\varepsilon m})$. On $|t|<\varepsilon$ put $t=\theta/a(m)$ and
$y=(j-m\mu)/a(m)$. The integrand becomes
$a(m)^{-1-r}(e^{-i\theta\mu/a(m)}\varphi(\theta/a(m)))^m(a(m)(e^{-i\theta/a(m)}-1))^re^{-i\theta y}$.
The first factor tends to $\exp(K_\alpha(-i\theta)^\alpha)$ by Lemma~\ref{lem:stable}(i) and the
second to $(-i\theta)^r$, with $|a(m)(e^{-i\theta/a(m)}-1)|\le|\theta|$. As in
Lemma~\ref{lem:trunc-local}(b), the modulus of the first factor is at most
$\exp(-c|\theta|^{\alpha-\varepsilon'})$ for $1\le|\theta|<\varepsilon a(m)$. Dominated
convergence, uniform in $y$ since $|e^{-i\theta y}|=1$, gives the limit
$(2\pi)^{-1}\int e^{K_\alpha(-i\theta)^\alpha}(-i\theta)^re^{-i\theta y}d\theta=g^{(r)}(y)$. For $q_m$
with $m\ge2$, $q_m(j)=\sum_{v\ge1}T(v)u_{m-1}(j-v)$, so the integrand carries the extra factor
$\hat\tau(t)=\sum_{v\ge1}T(v)e^{itv}$, which is continuous, bounded by $\mu$ and equal to $\mu$
at $0$. Since $a(m-1)/a(m)\to1$ and $g^{(r)}$ is uniformly continuous, the same proof applies.
The kernel $q_{m+1}$ is the same. The uniform bound follows from the domination for large $m$
and from $|\Delta^r\mathsf k_m|\le2^r\max(1,\mu)$ for the others.
\end{proof}

\subsection{Proof of Lemma~\ref{lem:diff-bigjump}}\label{app:diff-bigjump}

We write $f$, $T$, $\mu$, $\alpha$, $a$, $V_m$ for the plus sign as in the lemma, and
$u^{(w)}_m$ as in Lemma~\ref{lem:trunc-local}. By (LR), $f(L)\ge cT(L)/L$ for large $L$, so
by~\eqref{eq:potter-m}, with $r=L/a(m)$,
\begin{equation}\label{eq:target}
 m\,\frac{f(L)}L\ge c\,r^{-\alpha-2-\varepsilon'}a(m)^{-2},\qquad mT(\delta L)\le C_\delta r^{-\alpha+\varepsilon'}.
\end{equation}

\emph{The first difference.} For $m=1$ the claim is Lemma~\ref{lem:ld2}, since $L=n-\mu$ and
$f(n)/n\sim f(L)/L$. Let $m\ge2$. Fix $\delta\in(0,1/16]$ so small that
$(\alpha-\varepsilon')/(32\delta)>\alpha+2+\varepsilon'$, and put $w=\lceil\delta L\rceil$, so
$w\le2\delta L\le L/8$ for large $L$.

Step 1, no summand of size at least $w$. Put $m_1=\lfloor m/2\rfloor$, $m_2=m-m_1$,
$g_i=u^{(w)}_{m_i}$ and $Y=\lceil m_1\mu+L/2\rceil-1$. Summation by parts in the sum over $y\le Y$
gives the exact identity
\[
 \nabla u^{(w)}_m(n)=\sum_{y\le Y}\nabla g_1(y)\,g_2(n-y)+g_1(Y)\,g_2(n-1-Y)+\sum_{y>Y}g_1(y)\,\nabla g_2(n-y).
\]
In the first sum $n-y-m_2\mu\ge L/2$, and in the last $y-m_1\mu\ge L/2$. By
Lemma~\ref{lem:trunc-local}(b) and Lemma~\ref{lem:truncated} these two sums are at most
$Ca(m)^{-2}(mT(w))^{L/(4w)}$. In the middle term $Y-m_1\mu\ge L/4\ge2w$, so by
Lemma~\ref{lem:trunc-local}(a) and Lemma~\ref{lem:stable}(ii) it is at most
$Ca(m)^{-2}(mT(w))^{L/(16w)}$. Since $L/(16w)\ge1/(32\delta)$ and $mT(w)\to0$ as $r\to\infty$
by~\eqref{eq:target}, the whole is at most $Ca(m)^{-2}(C_\delta r^{-\alpha+\varepsilon'})^{1/(32\delta)}$,
which is $o(mf(L)/L)$ as $r\to\infty$ by~\eqref{eq:target} and the choice of $\delta$.

Step 2, a giant summand. Let $I$ be the first index with $X^\plus_I\ge w$, and let $h_i(z)$ be
the probability that $X^\plus_l<w$ for $l<i$ and that the sum of the $X^\plus_l$ with $l\le m$,
$l\ne i$, equals $z$. Then
\[
 \nabla\Prob(S^\plus_m=\cdot,\,I=i)(n)=\sum_{z\le n-1-w}h_i(z)\nabla f(n-z)-h_i(n-w)f(w).
\]
Fix $\varepsilon_1>0$. On the main range $|z-(m-1)\mu|\le\varepsilon_1L$ we have
$n-z=L+\mu-(z-(m-1)\mu)$, so Lemma~\ref{lem:ld2} and Lemma~\ref{lem:rv}(i) give
$\nabla f(n-z)=(\alpha+1)(f(L)/L)(1+\theta_1)$ with $|\theta_1|\le\omega(\varepsilon_1)+o(1)$ as
$L\to\infty$, where $\omega(\varepsilon_1)\to0$. The mass of the main range lies between
$(1-T(w))^{i-1}-p_m$ and $(1-T(w))^{i-1}$, with $p_m=\Prob(|S^\plus_{m-1}-(m-1)\mu|>\varepsilon_1L)$.
Off the main range $n-z\ge w\ge\delta L$, so $|\nabla f(n-z)|\le C_\delta f(L)/L$ by (LD) and
Lemma~\ref{lem:rv}(ii), and the mass is at most $p_m$. By~\eqref{eq:right-tail}
and Lemma~\ref{lem:stable}(iv),
$p_m\le CmT(\varepsilon_1L)+e^{C_-}e^{-\varepsilon_1L/a(m)}$, which tends to $0$ as $r\to\infty$,
uniformly in $m$, by~\eqref{eq:potter-m}. Next,
$\sum_{i\le m}(1-T(w))^{i-1}=m(1+O(mT(w)))$. Finally the deviation of $n-w$ from $(m-1)\mu$ is at
least $L/2$, so $h_i(n-w)\le u^\plus_{m-1}(n-w)\le CmT(L)/a(m)$ by Lemma~\ref{lem:stable}(iii),
and $f(w)\le C_\delta f(L)$. Summed over $i\le m$ these terms are at most $Cm^2T(L)f(L)/a(m)$,
which relative to $mf(L)/L$ is $CmT(L)L/a(m)\le Cr^{1-\alpha+\varepsilon'}$. Adding Steps 1 and 2,
$\nabla V_m(n)=(\alpha+1)mf(L)/L\,(1+\theta_2)$ with $|\theta_2|\le C\omega(\varepsilon_1)+o(1)$ as
$r\to\infty$, uniformly in $m$. Choose $\varepsilon_1$ and then $K$.

\emph{The bound on $V_m(n)$.} The same decomposition without differences gives
\[
 V_m(n)\le u^{(w)}_m(n)+m\sup_{y\ge w}f(y).
\]
By Lemma~\ref{lem:trunc-local}(a) the first term is at most $Ca(m)^{-1}(mT(w))^{L/(4w)}$, which
is $o(mf(L))$ by the same comparison. The second is at most $C_\delta mf(L)$.

\emph{The difference in the index.} We have $\delta V_m(n)=\sum_xf(x)[V_m(n-x)-V_m(n)]$. Fix
$\varepsilon_1,\varepsilon_2>0$. For $x\le\varepsilon_1L$ the bracket is
$\sum_{i<x}\nabla V_m(n-i)$, and the first bound at the deviations $L-i\ge(1-\varepsilon_1)L$,
with Lemma~\ref{lem:rv}(i), makes it $x(\alpha+1)mf(L)/L$ up to a factor $1+\theta$ with
$|\theta|\le\varepsilon+\omega(\varepsilon_1)$. Since $\sum_{x\le\varepsilon_1L}xf(x)\to\mu$, this
part is $\mu(\alpha+1)mf(L)/L$ up to the same factor and $o(1)$. The term
$-V_m(n)\sum_{x>\varepsilon_1L}f(x)$ is at most $Cmf(L)T(\varepsilon_1L)=o(mf(L)/L)$. The rest is
$\E[f(n-S^\plus_m);\,n-S^\plus_m>\varepsilon_1L]$. On $|S^\plus_m-m\mu|\le\varepsilon_2L$ the
integrand is $f(L)(1+O(\omega(\varepsilon_2)))$, and this event has probability $1-o(1)$ as
$r\to\infty$. On $S^\plus_m-m\mu>\varepsilon_2L$ the integrand is at most $C_{\varepsilon_1}f(L)$,
with probability at most $CmT(\varepsilon_2L)$. On $S^\plus_m-m\mu<-\varepsilon_2L$ it is at most
$Cf(L)$, with probability at most $e^{C_-}e^{-\varepsilon_2r}$. So the rest is $f(L)(1+o(1))$ up
to $O(\omega(\varepsilon_2))f(L)$.

\emph{The difference of the $q$ kernels.} We have
$q_{m+2}(n)-q_{m+1}(n)=\sum_xf(x)[q_{m+1}(n-x)-q_{m+1}(n)]$, with $q_{m+1}=0$ at nonpositive
arguments. For $x\le\varepsilon_1L$ the bracket is $\sum_{l=1}^x\delta V_m(n-l)$, and the previous
step makes this part $\mu[f(L)+\mu(\alpha+1)mf(L)/L]$ up to small relative errors. For $m=0$ use
$\delta V_0(y)=f(y)$ instead. The term $-q_{m+1}(n)\sum_{x>\varepsilon_1L}f(x)$ is at most
$C(T(L)+mf(L))T(\varepsilon_1L)$, which is $o(f(L)+mf(L)/L)$. The rest equals
$\E\sum_{y=S^\plus_m+1}^{S^\plus_{m+1}}f(n-y)\ones_{\{n-y>\varepsilon_1L\}}$. On the event that
$|S^\plus_m-m\mu|\le\varepsilon_2L$ and $X^\plus_{m+1}\le\varepsilon_2L$ the sum is
$X^\plus_{m+1}f(L)(1+O(\omega(2\varepsilon_2)))$, and the expectation of $X^\plus_{m+1}$ on this
event tends to $\mu$, by independence. On the complement the sum is at most
$C_{\varepsilon_1}f(L)X^\plus_{m+1}$, whose expectation there is at most
$C_{\varepsilon_1}f(L)(\E[X_\plus;X_\plus>\varepsilon_2L]+\mu\Prob(|S^\plus_m-m\mu|>\varepsilon_2L))=o(f(L))$.
Adding the parts gives the claim.

\emph{The bound on $q_{m+1}$.} Split $q_{m+1}(n)=\sum_xV_m(x)T(n-x)$ at $x=n-L/2$. The lower part
is at most $T(L/2)\le CT(L)$. In the upper part the deviation of $x$ is at least $L/2$, so the
bound on $V_m$ gives $V_m(x)\le Cmf(L)$, and the weights $T(n-x)$ add up to at most $\mu$.
\qed

\subsection{The big-jump zone}\label{app:zone}

From here on we assume $\alpha_\plus=\alpha_\minus=\alpha$, fix $\theta_0>0$, and consider only
the $N$ with $\min_\sgn\theta_{\sgn,N}\ge\theta_0$. ``For all large $N$'' means for all large
such $N$, and constants may depend on $\theta_0$. Under these conditions $a_\plus(m)$ and
$a_\minus(m)$ are comparable, uniformly in $m\in[\varepsilon N,2N]$ for every fixed
$\varepsilon>0$. Indeed $a_\plus/a_\minus$ is slowly varying, its value at $N/s$ lies in
$[\theta_0,1/\theta_0]$, and Lemma~\ref{lem:rv}(i) transfers this to $m\in[\varepsilon N,2N]$.

Put $L=L_{N-j}=N-js/\mu_\minus$, $i_0=j/\mu_\minus$ and
\[
 \mathcal D_N(j)=\frac{2s}{\mu_\minus^2}f_\plus(L)+\frac{js^2}{\mu_\minus^3}\,(\alpha+1)\,\frac{f_\plus(L)}L .
\]
For a plus start this is $(1+o(1))$ times the derivative in $j$ of the kernel~\eqref{eq:G}, by
$\gamma(j)\sim js/\mu_\minus^2$ and $u_\minus(j)\to1/\mu_\minus$.

\begin{theorem}[a difference form of Theorem~\ref{thm:bigjump}(b)]\label{thm:zone-diff}
Assume (SPAN) and (LR) for both signs and (LD2) for the plus sign, and the conditions above. Let
the start be fresh with either sign or fair, and fix $\eta\in(0,1-1/\alpha)$. Then
\[
 \lim_{K\to\infty}\ \limsup_{N\to\infty}\ \sup_{N^{1-\eta}\le j\le(1-p)N-KB_N}
 \Bigl|\frac{D_N(j)}{\mathcal D_N(j)}-1\Bigr|=0 .
\]
\end{theorem}

\begin{proof}
Take the plus start and write $n=N-j$. In the range of the theorem $L\ge KB_N$. The summations
by parts in the proof of Theorem~\ref{thm:first-diff} are exact for every $j\ge1$, and give
$D_N=\mathrm I+\mathrm{II}+\mathrm{III}+\mathrm{IV}$ with
\begin{align*}
 \mathrm I&=Q_1(n-1)\Delta u(j)+\sum_{i\ge1}\delta Q_i(n-1)\Delta\bar A_i(j),&
 \mathrm{II}&=\sum_mu^\minus_m(j)\,\delta V_m(n-1),\\
 \mathrm{III}&=\sum_{i\ge1}\delta V_i(n-1)\,u^\minus_i(j),&
 \mathrm{IV}&=\sum_m\nabla V_m(n)\,q^\minus_m(j),
\end{align*}
where $u=u_\minus$, $Q_m=q_{m+1}$ and $\delta Q_i=q_{i+2}-q_{i+1}$. Only $m\le j+1$ occurs. Put
$L_m=n-1-m\mu_\plus$, so $L_{i_0}=L-1$. Fix $M\ge2$, then $K$, then $N$, and write
$I=I^\minus_j(M)$ for the index window of Lemma~\ref{lem:index}.

The window. For $m\in I$, $|m-i_0|\le Ma_\minus(j/\mu_\minus)/\mu_\minus\le CMB_N$, so
$L_m=L(1+O(M/K))$. Also $m=i_0(1+o(1))$, since $a_\minus(j)=o(j)$ and $j\ge N^{1-\eta}$. Since
$a_\plus(m)\le a_\plus(N)\le CB_N$, we have $L_m\ge Ka_\plus(m)/C$. So
Lemma~\ref{lem:diff-bigjump} and Lemma~\ref{lem:rv}(i) give, uniformly in $m\in I$,
\[
 \nabla V_m(n)=(\alpha+1)m\frac{f(L)}L(1+o(1)),\quad
 \delta V_m(n-1)=\Bigl[f(L)+\mu_\plus(\alpha+1)i_0\frac{f(L)}L\Bigr](1+o(1)),
\]
and $\delta Q_m(n-1)=[2\mu_\plus f(L)+\mu_\plus^2(\alpha+1)i_0f(L)/L](1+o(1))$, where $f=f_\plus$ and
$o(1)$ tends to $0$ as $N\to\infty$ and then $K\to\infty$.

The other indices with $m\le m_*=i_0+L/(2\mu_\plus)$. For them $L_m\ge L/3$, so
Lemma~\ref{lem:diff-bigjump} gives $|\nabla V_m(n)|\le Cmf(L)/L$ and $|\delta V_m(n-1)|$,
$|\delta Q_m(n-1)|\le C(f(L)+i_0f(L)/L)$, using $m\le j+1\le Ci_0$. By Lemma~\ref{lem:index}(a)
the masses $\sum_{m\notin I}u^\minus_m(j)$ and $\sum_{m\notin I}q^\minus_m(j)$ are at most
$\eta_\minus(M)+o(1)$. Moreover $\sum_{m\notin I}mq^\minus_m(j)\le i_0(2\eta_\minus(M)+o(1))$. For
$m\le2i_0$ this follows from the masses, and $\sum_{m>2i_0}mq^\minus_m(j)\le(j+1)\Prob(J_\minus(j)>2i_0)$
is at most $(j+1)e^{C_-}e^{-cj/a_\minus(j)}=o(1)$ by Lemma~\ref{lem:stable}(iv).

The bad indices $m_*<m\le j+1$. For them $m\mu_\minus-j\ge\mu_\minus L/(2\mu_\plus)=c_*L$, a left
deviation of the minus runs. By~\eqref{eq:left-local}, $u^\minus_m(x)\le Ca_\minus(m)^{-1}e^{-c_*L/(2a_\minus(m))}$
for $x\le j$, and $q^\minus_m(j)=\sum_vT_\minus(v)u^\minus_{m-1}(j-v)$ obeys the same bound with
$c_*L/2$ in place of $c_*L$. For the plus factors we use the crude bounds of
Lemma~\ref{lem:fourier-diff}, $|\nabla V_m|\le Ca_\plus(m)^{-2}$,
$|\delta V_m|\le\mu_\plus\sup|\nabla V_m|$ and $|\delta Q_m|\le\mu_\plus\sup|\nabla q_{m+1}|\le Ca_\plus(m)^{-2}$.
Bad indices exist only if $m_*<j+1$, which forces $L\le Cj$ and so $j\ge cN$. Then
$m\in[cN,N]$, and $a_\plus(m)$, $a_\minus(m)$ are comparable to $A=a_\minus(j)$. Summing the
exponential bound over $m$ gives at most $Ce^{-cL/A}$ for the minus masses, so every bad
contribution is at most $CA^{-2}e^{-cL/A}$. With $\rho=L/A\ge cK$ we have
$\mathcal D_N(j)\ge ci_0f(L)/L\ge c\rho^{-\alpha-2-\varepsilon'}A^{-2}$ by~\eqref{eq:potter-m},
since $i_0T_\plus(A)$ is bounded below. So the bad contributions are at most
$C\rho^{\alpha+2+\varepsilon'}e^{-c\rho}\mathcal D_N(j)$, which is small as $K\to\infty$.

The four parts. By the window, the masses and $\sum_mmq^\minus_m(j)=\E J_\minus(j)=\sum_{x<j}u(x)=i_0(1+o(1))$,
\begin{align*}
 \mathrm{IV}&=(\alpha+1)\,i_0\,\frac{f(L)}L\bigl(1+O(\eta_\minus(M))+o(1)\bigr),\\
 \mathrm{II},\ \mathrm{III}&=\frac1{\mu_\minus}\Bigl[f(L)+\mu_\plus(\alpha+1)i_0\frac{f(L)}L\Bigr]\bigl(1+O(\eta_\minus(M))+o(1)\bigr),
\end{align*}
where we used $\sum_{m\in I}u^\minus_m(j)=u(j)+O(\eta_\minus(M))+o(1)$ and $u(j)\to1/\mu_\minus$.
In I, $Q_1(n-1)=q_2(n-1)\le CT_\plus(n)\le CT_\plus(L)$ and $|\Delta u(j)|\le CT_\minus(j)$, so the
first term is at most $CLT_\minus(j)f(L)\le CNT_\minus(N^{1-\eta})f(L)=o(f(L))$, since
$(1-\eta)\alpha>1$. For the sum let $i_*=\lfloor i_0\rfloor\in I$ and write it as
$\delta Q_{i_*}(n-1)\Delta(v-u)(j)+\sum_i(\delta Q_i-\delta Q_{i_*})(n-1)\Delta\bar A_i(j)$, using
$\sum_i\bar A_i=v-u$. Here $\Delta(v-u)(j)\to1/\mu_\minus^2$ by the proof of
Theorem~\ref{thm:first-diff}. By fact (2) of that proof,
$\Delta\bar A_i(j)=\sum_{x<j}u^\minus_i(x)d_{j+1-x}+u^\minus_i(j)u(1)$. Fix $R'$. The terms with
$j-x>R'$ contribute at most $(\sum_{r>R'}|d_r|)\sup_{x\le j}\sum_iu^\minus_i(x)(|\delta Q_i|+|\delta Q_{i_*}|)$,
and the index classes at the level $x$ bound the last sum by $C(f(L)+i_0f(L)/L)+o(\mathcal D_N(j))$.
For the terms with $j-x\le R'$, the indices $i\in I^\minus_x(M)$ satisfy
$|\delta Q_i-\delta Q_{i_*}|=o(f(L)+i_0f(L)/L)$ by the window, since $|x-j|\le R'$, and the other
indices carry mass at most $\eta_\minus(M)+o(1)$. Hence
\[
 \mathrm I=\frac1{\mu_\minus^2}\Bigl[2\mu_\plus f(L)+\mu_\plus^2(\alpha+1)i_0\frac{f(L)}L\Bigr]
 +O\Bigl(\eta_\minus(M)+\sum_{r>R'}|d_r|+o(1)\Bigr)\Bigl(f(L)+i_0\frac{f(L)}L\Bigr).
\]
Adding the four parts, the coefficient of $f(L)$ is $2\mu_\plus/\mu_\minus^2+2/\mu_\minus=2s/\mu_\minus^2$
and that of $(\alpha+1)i_0f(L)/L$ is $(\mu_\plus+\mu_\minus)^2/\mu_\minus^2=s^2/\mu_\minus^2$. So
$D_N(j)=\mathcal D_N(j)(1+O(\eta_\minus(M)+\sum_{r>R'}|d_r|)+o(1))$, since
$\mathcal D_N(j)\ge c(f(L)+i_0f(L)/L)$. Let $N\to\infty$, then $K\to\infty$, then $M,R'\to\infty$.

For a minus start the indices shift by one. The extra term $V_1(n-1)f_\minus(j)$ in III is
$O(f_\plus(L)f_\minus(j))=o(f(L))$, and $\sum_mmq^\minus_{m+1}(j)=\E J_\minus(j)-1$. The leading
terms are the same, and the fair start follows by linearity.
\end{proof}

\begin{proof}[Proof of Theorem~\ref{thm:mode-count}(iv)]
Choose $K_0$ such that $|D_N/\mathcal D_N-1|\le\frac12$ on $N^{1-\eta}\le j\le(1-p)N-K_0B_N$ for all
large $N$. Then $D_N>0$ there. On $(x_*+\eta')\Lambda_N\le j\le N^{1-\eta}$ we have $D_N>0$ by
Theorem~\ref{thm:ends}(c). At $j=N^{1-\eta}$ both statements hold, and
$\mathcal D_N\sim c_{\mathrm a}T_\plus(N)/N$ matches the limit in Theorem~\ref{thm:first-diff}.
\end{proof}

The proof uses the comparability of the scales only for the bad indices, which occur only for
$j\ge cN$. Without it the same bound needs $e^{-cL/a_\minus(j)}$ to beat
$(a_\minus(j)/a_\plus(j))^2$, which fails near the window of Theorem~\ref{thm:window}.

\subsection{The bulk}\label{app:bulk}

We keep the conditions of Section~\ref{app:zone}. Fix a sign and drop it, as in
Section~\ref{app:fourier}, put $A=a(j/\mu)$ and let $I_j(K)$ be the index window of
Lemma~\ref{lem:index}.

\begin{lemma}[a pointwise bound]\label{lem:pointwise}
Assume (RV) and (SPAN), and let $r\in\{1,2\}$. For $m\ge1$ and integers $x$ put $z=x-m\mu$. Then
\begin{align*}
 |\Delta^ru_m(x)|&\le Ca(m)^{-1-r}\min\{1,mT(z/4)\}&&(z\ge0),\\
 |\Delta^ru_m(x)|&\le Ca(m)^{-1-r}e^{-|z|/(8a(m))}&&(z<0),
\end{align*}
and $\sum_x|\Delta u_m(x)|\le C/a(m)$.
\end{lemma}

\begin{proof}
The bound with $\min\{1,\cdot\}=1$ is Lemma~\ref{lem:fourier-diff}. So for $z\ge0$ we may assume
$mT(z/4)\le1$ and $z\ge24$, and for $z<0$ we may assume $|z|\ge24\max(1,a(m))$. For $m\le2$ the
claims are direct, since $|\Delta^ru_1(x)|\le4T(x)$, $u_2(y)\le2T(y/2)$, and $x\ge m$ gives
$|z|\le2\mu$ when $z<0$.

Let $r=1$, $m\ge2$ and $z\ge24$. Write $u_m=u_{(1)}*u_{(2)}$ with block sizes $m_1=\lfloor m/2\rfloor$
and $m_2=m-m_1$, means $\nu_i=m_i\mu$, and put $c=\lceil\nu_1+z/2\rceil$. Then
$\Delta u_m(x)=\sum_{x_1}u_{(1)}(x_1)\Delta u_{(2)}(x-x_1)$. The part $x_1\ge c-1$ is at most
$\sup|\Delta u_{(2)}|\Prob(S_{m_1}-\nu_1\ge z/2-1)\le Ca(m)^{-2}mT(z/3)$, by
Lemma~\ref{lem:fourier-diff} and~\eqref{eq:right-tail}. On the part $x_1\le c-2$ summation by parts
gives
\[
 -\sum_{x_1\le c-2}\Delta u_{(1)}(x_1-1)\,u_{(2)}(x-x_1)+u_{(1)}(c-2)\,u_{(2)}(x-c+1).
\]
The sum is at most $\sup|\Delta u_{(1)}|\Prob(S_{m_2}-\nu_2\ge z/2)$, and in the last term
$u_{(1)}(c-2)\le Cm_1T(z/3)/a(m)$ by Lemma~\ref{lem:stable}(iii) and $u_{(2)}\le C/a(m)$.

Let $r=2$, $m\ge3$ and $z\ge24$. Use three blocks $u_{(1)},u_{(2)},u_{(3)}$ of sizes
$\lfloor m/3\rfloor$ or $\lceil m/3\rceil$, with means $\nu_i$, and cuts
$c_i=\lceil\nu_i+z/3\rceil$ for $i=1,2$. Split $u_m(x)$ into the parts with $x_1\ge c_1$, with
$x_1<c_1$ and $x_2\ge c_2$, and with $x_1<c_1$ and $x_2<c_2$, where $x_i$ is the value of block $i$.
The second difference of the first part is at most
$\Prob(S_{(1)}-\nu_1\ge z/3)\sup|\Delta^2u_{(2)}*u_{(3)}|\le CmT(z/4)a(m)^{-3}$, and the second
part is the same. In the third part block $3$ lies above $\nu_3+z/3-2$. Put
$H(y)=\sum_{x_2<c_2}u_{(2)}(x_2)u_{(3)}(y-x_2)$ and $F(x)=\sum_{x_1<c_1}u_{(1)}(x_1)H(x-x_1)$.
Summation by parts as for $r=1$ gives
\[
 \Delta^2F(x)=-\sum_{x_1<c_1}\Delta u_{(1)}(x_1-1)\,\Delta H(x-x_1)+u_{(1)}(c_1-1)\,\Delta H(x-c_1),
\]
\[
 \Delta H(y)=-\sum_{x_2<c_2}\Delta u_{(2)}(x_2-1)\,u_{(3)}(y-x_2)+u_{(2)}(c_2-1)\,u_{(3)}(y-c_2).
\]
Expanding, each term is a product of factors of three kinds. Either it contains
$\|\Delta u_{(1)}\|_1\sup|\Delta u_{(2)}|$ and a tail probability of block $3$ at deviation at
least $z/3$, or it contains a local probability $u_{(i)}(c_i-1)$ at deviation at least $z/3-1$,
together with $\|\Delta u_{(1)}\|_1$ or $\|\Delta u_{(2)}\|_1$ or $u_{(1)}(c_1-1)$, and
$\sup u_{(3)}$. By~\eqref{eq:right-tail}, Lemma~\ref{lem:stable}(ii) and~(iii) and the bound
$\|\Delta u_{(i)}\|_1\le C/a(m)$ proved below from the case $r=1$, each is at most
$Ca(m)^{-3}mT(z/4)$.

For $z<0$ mirror the argument. The cuts become $\lfloor\nu_i+z/2\rfloor$ or $\lfloor\nu_i+z/3\rfloor$,
the summations by parts run over the sums restricted from below, and~\eqref{eq:left-local} and
Lemma~\ref{lem:stable}(iv) replace Lemma~\ref{lem:stable}(iii) and~\eqref{eq:right-tail}. The
deviations are at least $|z|/4$, which gives the exponent.

For the total variation, the two-block argument for $r=1$, mirrored when $z<0$, gives for every
$x$
\[
 |\Delta u_m(x)|\le Ca(m)^{-2}\sum_{i=1,2}\Prob\bigl(|S_{m_i}-\nu_i|\ge|z|/2-3\bigr)+u_{(1)}(c_x)\,u_{(2)}(x-c_x+e_x),
\]
with $|c_x-\nu_1-z/2|\le3$ and $e_x\in\{-1,1\}$. Summing over $x$, the probabilities add up to at
most $C\E|S_{m_i}-\nu_i|+C\le Ca(m)$. The point $c_x$ takes each value for at most three $x$, so
the products add up to at most $3\sup u_{(2)}\le C/a(m)$.
\end{proof}

\begin{lemma}[second differences in the big-jump range]\label{lem:second-bigjump}
Assume (SPAN), (LR) and (LD2). There are $K_1$ and $L_0$ such that for all $m\ge1$ and all $n$
with $L=n-m\mu\ge\max(K_1a(m),L_0,m\mu)$,
\[
 |\nabla^2u_m(n)|\le Cm\frac{f(L)}{L^2},\qquad |\nabla u_m(n)|\le Cm\frac{f(L)}L .
\]
\end{lemma}

\begin{proof}
The first-difference bound is Lemma~\ref{lem:diff-bigjump}. For $m=1$ the second is (LD2). Let
$m\ge2$, fix $\delta>0$ small and put $w=\lceil\delta L\rceil$. The condition $L\ge m\mu$ gives
$mT(\delta L)\le C_\delta LT(L)\le C_\delta L^{1-\alpha+\varepsilon'}$.

No summand of size at least $w$. With $g_i$ as in the proof of Lemma~\ref{lem:diff-bigjump},
$\nabla^2u^{(w)}_m=(\nabla g_1)*(\nabla g_2)$ exactly. For each $y$ one of the blocks exceeds its
mean by at least $L/2$. Since $\sum_{y\ge Y}|\nabla g(y)|\le2\Prob(S\ge Y-1)$ for the law $g$ of a
sum $S$, Lemma~\ref{lem:trunc-local}(b) and Lemma~\ref{lem:truncated} bound this part by
$Ca(m)^{-2}(mT(w))^{L/(8w)}\le CL^{-(\alpha-1-\varepsilon')/(16\delta)}$. This is $o(f(L)/L^2)$ for
$\delta$ small, since $f(L)/L^2\ge L^{-\alpha-3-\varepsilon'}$ for large $L$.

A giant summand. With $I$ and $h_i$ as before, the second difference of $\Prob(S_m=\cdot,I=i)$
at $n$ is the sum of three terms, as in the proof of Lemma~\ref{lem:plus-second}(i). The main
sum is at most $C_\delta f(L)/L^2$ by (LD2), since $n-z\ge\delta L$. In the first boundary term
$h_i(n-1-w)\le u_{m-1}(n-1-w)\le CmT(L)/a(m)$ by Lemma~\ref{lem:stable}(iii) and
$|f(w+1)-f(w)|\le C_\delta f(L)/L$. Summed over $i$ this is at most $Cmf(L)/L^2$ times
$CmT(L)L/a(m)\le C$, by~\eqref{eq:potter-m}. For the second boundary term we apply the proof of
the first bound of Lemma~\ref{lem:diff-bigjump}, as an upper bound, to the law $h_i$ at the level
$n-w$, with the threshold $w$ itself. The deviation is at least $L/2$, a truncated variable is
never a giant summand, and the no-giant part is bounded as above. This gives
$|h_i(n-w)-h_i(n-1-w)|\le Cmf(L)/L$, so the term is at most $C_\delta f(L)\,mf(L)/L$. Summed over
$i$ it is $Cmf(L)/L^2$ times $CmLf(L)\le CmT(L)\le C$.
\end{proof}

\begin{lemma}[tails of the kernel differences]\label{lem:tails}
Assume (SPAN), (LR), (LD2), let $\mathsf k$ be one of $u_m$, $q_m$, $q_{m+1}$, and let
$r\in\{1,2\}$.
\begin{enumerate}[(a)]
\item For $K\ge2$ and $j\to\infty$, $\sum_{m\notin I_j(K)}|\Delta^r\mathsf k_m(j)|\le C(K^{1-\alpha+\delta}+e^{-cK})A^{-r}+o(A^{-r})$,
and $\sum_m|\Delta^r\mathsf k_m(i)|\le Ca(i)^{-r}$ for all $i\ge1$.
\item Fix $\varepsilon>0$ and put $m_\varepsilon=\varepsilon N/(8\mu_{\max})$. Uniformly in levels
$x\in[\varepsilon N,N]$, $|\Delta^r\mathsf k_m(x)|\le Ca(N)^{-1-r}$ for $m\ge m_\varepsilon$. For
$m<m_\varepsilon$, $|\Delta^ru_m(x)|\le Cmf(x)x^{-r}$, and $|\Delta^rq_m(x)|$ and
$|\Delta^rq_{m+1}(x)|$ are at most $Cf(x)x^{1-r}(1+m/x)$.
\item Let $\mathsf k'$ be a kernel of the other sign and let $x'\in[\varepsilon N,N]$. Then
\[
 \sum_{m<m_\varepsilon}|\Delta^r\mathsf k_m(x)|\,\mathsf k'_m(x')=o(B_N^{-1-r}),
\]
and $\sum_{m<m_\varepsilon}|\Delta\mathsf k_m(x)|\,|\Delta\mathsf k'_m(x')|=o(B_N^{-3})$ when one
of $\mathsf k,\mathsf k'$ is a $u$ kernel and the other a $q$ kernel.
\end{enumerate}
\end{lemma}

\begin{proof}
(a) for $u$. Split $m\notin I_j(K)$ into three ranges. For $m\le j/(2\mu)$ we have
$L=j-m\mu\ge j/2\ge\max(K_1a(m),L_0,m\mu)$ for large $j$, so Lemma~\ref{lem:second-bigjump} gives
$|\Delta^ru_m(j)|\le Cmf(j)j^{-r}$, and the range contributes at most $Cj^{2-r}f(j)\le CT(j)j^{1-r}$.
This is $o(A^{-r})$, since $T(j)a(j)\to0$ and $T(j)a(j)^2/j\to0$ for $\alpha>1$. For
$j/(2\mu)<m<(j-KA)/\mu$ we have $a(m)\asymp A$, and Lemma~\ref{lem:pointwise} bounds the terms by
$CA^{-1-r}jT(z/4)$ with $z=j-m\mu\ge KA$ running through a progression of step $\mu$. By
Karamata's theorem the sum is at most $CA^{-1-r}j\,KA\,T(KA)\le CK^{1-\alpha+\delta}A^{-r}$, since
$jT(A)\to\mu$. For $m>(j+KA)/\mu$ only $m\le j+2$ contributes, and Lemma~\ref{lem:pointwise}
gives at most $CA^{-1-r}\sum_{l\ge0}e^{-(KA+l\mu-2)/(CA)}\le CA^{-r}e^{-cK}$. The window $I_j(2)$
contributes at most $CA^{-r}$, which gives the second claim.

(a) for $q$. For $m=1$, $\Delta^rq_1(j)=\Delta^rT(j)=O(f(j))=o(A^{-r})$. For $m\ge2$,
$\Delta^rq_m(j)=\sum_{v\ge1}T(v)\Delta^ru_{m-1}(j-v)$. Fix $v_0$. For $v\le v_0$ and large $j$,
$m\notin I_j(K)$ implies $m-1\notin I_{j-v}(K/2)$, so these shifts contribute at most $\mu$ times
the bound for $u$ with $K/2$. The shifts $v_0<v\le j/2$ contribute at most
$(\sum_{v>v_0}T(v))CA^{-r}$, by the second claim for $u$. For the shifts $v>j/2$ we do not take
absolute values inside the sum over $v$. For each $m'=m-1$ we sum by parts $r$ times in
$i=j-v<j/2$, which moves the differences onto $T(j-i)$. By (LD),
$|\Delta^r_iT(j-i)|\le Cf(j)j^{1-r}$ there, and $\sum_{m'}u_{m'}(i)\le1$. The boundary terms are
$T(j/2+O(1))|\Delta^{r-1}u_{m'}(\lfloor j/2\rfloor+O(1))|$ and, for $r=2$, a term
$f(j/2+O(1))u_{m'}(\lfloor j/2\rfloor+O(1))$. Summed over all $m'$, by the second claim for $u$,
these shifts contribute at most $CT(j)j^{1-r}+CT(j)a(j)^{1-r}+Cf(j)=o(A^{-r})$. Let $v_0\to\infty$.
The kernel $q_{m+1}=u_m*\tau$ is the same.

(b) For $m\ge m_\varepsilon$ this is Lemma~\ref{lem:fourier-diff}, since $a(m)\ge a(m_\varepsilon)\ge ca(N)$,
and $|\Delta^rq_m|\le\mu\sup|\Delta^ru_{m-1}|$. For $m<m_\varepsilon$ and $u$ we have
$x-m\mu\ge x/2\ge\max(K_1a(m),L_0,m\mu)$ for large $N$, and Lemma~\ref{lem:second-bigjump} applies.
For $q_m$ split $\sum_vT(v)\Delta^ru_{m-1}(x-v)$ at $v=x/2$. The part $v\le x/2$ is at most
$C\mu mf(x)x^{-r}$ by the bound for $u$. For the part $v>x/2$ the summation by parts above,
for the single index $m-1$, gives at most $Cf(x)x^{1-r}$ plus boundary terms. These are
$T(x/2)u_{m-1}(x/2)\le CT(x)mf(x)\le Cf(x)$ for $r=1$, and
$T(x/2)|\Delta u_{m-1}(x/2)|+f(x/2)u_{m-1}(x/2)\le Cf(x)x^{-1}(1+m/x)$ for $r=2$, using
$mT(x)\le xT(x)\to0$ and $xT(x)\le1$.

(c) Write $T'$, $f'$ and $a'$ for the other sign. By Lemma~\ref{lem:small},
$\mathsf k'_m(x')\le C(1+m/a'(m))T'(N)$ for $m<m_\varepsilon$. By (b) and $f(x)\le Cf(N)$, the
first sum is at most
\begin{align*}
 &C\frac{f(N)}{N^r}T'(N)\sum_{m<\varepsilon N}m\Bigl(1+\frac m{a'(m)}\Bigr)
 +Cf(N)N^{1-r}T'(N)\sum_{m<\varepsilon N}\Bigl(1+\frac m{a'(m)}\Bigr)\\
 &\qquad\le C\frac{N^{2-r}T(N)T'(N)}{a'(N)}.
\end{align*}
Both signs have the index $\alpha$, so the right side is regularly varying with the index
$2-r-2\alpha-1/\alpha$, and $B_N^{-1-r}$ with the index $-(1+r)/\alpha$. The claim follows from
$1-2\alpha-1/\alpha<-2/\alpha$ and $-2\alpha-1/\alpha<-3/\alpha$, which hold for $\alpha>1$. For
the mixed sum, (b) bounds the product of a $u$ difference and a $q$ difference by
$Cmf(N)N^{-1}\cdot Cf'(N)$, so the sum is at most $CNf(N)f'(N)\le CT(N)T'(N)/N$, whose index
$-2\alpha-1$ is below $-3/\alpha$.
\end{proof}

\begin{lemma}[the bell shape]\label{lem:bell}
Let $\alpha_\plus=\alpha_\minus=\alpha$ and $\theta\in\Theta$ with $\min_\sgn\theta_\sgn>0$. Then
$h_\theta$ extends to an entire function, $h'_\theta$ has exactly one zero $x_\theta$, and
$h'_\theta>0$ on $(-\infty,x_\theta)$, $h'_\theta<0$ on $(x_\theta,\infty)$ and
$h''_\theta(x_\theta)<0$. On $\Theta_{\theta_0}=\{\theta\in\Theta:\min_\sgn\theta_\sgn\ge\theta_0\}$
the map $\theta\mapsto x_\theta$ is continuous, and there are $\varepsilon_0,c_0>0$ with
$h''_\theta\le-c_0$ on $|x-x_\theta|\le\varepsilon_0$. For every $R$ and $\varepsilon>0$ there
is $c>0$ with $|h'_\theta(x)|\ge c$ for $|x|\le R$, $|x-x_\theta|\ge\varepsilon$ and
$\theta\in\Theta_{\theta_0}$.
\end{lemma}

\begin{proof}
By the proof of Lemma~\ref{lem:profile}, $W_\theta$ is the sum of two independent strictly
$\alpha$-stable variables, so it is strictly $\alpha$-stable, and it is nondegenerate. Its
characteristic function satisfies $|\widehat h_\theta(t)|=e^{-c_\theta|t|^\alpha}$ with
$c_\theta>0$, so $\int e^{|t||\operatorname{Im}z|}|\widehat h_\theta(t)|\,dt<\infty$ for every
$z\in\mathbb C$, and $h_\theta$ is entire. By Kwa\'snicki~\cite[Corollary~1.3]{kwasnicki2020},
every stable density on $\R$ is bell-shaped in the sense of his Definition~2.1, that is, for every
$n\ge0$ the derivative $h^{(n)}_\theta$ tends to $0$ at $\pm\infty$ and changes sign exactly $n$
times.

Let $g=h^{(n)}_\theta$. It is analytic and not identically zero, so its zeros are isolated and of
finite order. If $g$ has $D$ consecutive zeros, then $g'$ changes sign in each of the $D-1$ gaps
between them, since on a gap $|g|$ rises from $0$ and returns to $0$. As $g'$ changes sign
exactly $n+1$ times, $g$ has finitely many zeros, say $D$ of them. Of these, $S$ have odd order
and $T_0$ even order, and $S=n$ since $g$ changes sign exactly $n$ times. The derivative $g'$
changes sign once in each gap, once in each of the two unbounded intervals, where $g\ne0$ and
$g\to0$ at infinity, and once at each zero of even order, where $g'$ has a zero of odd order.
These are distinct places, so $n+1\ge D+1+T_0=n+1+2T_0$, and $T_0=0$. This holds for every $n\ge1$.
If a zero of $g$ had order $k\ge3$, it would be a zero of $g'$ of even order. So all zeros of
every derivative are simple. For $n=1$, $h'_\theta$ has one simple zero, $h_\theta>0$ tends to
$0$ at $\pm\infty$, so the zero is the maximum, $h'_\theta$ has the stated signs and
$h''_\theta(x_\theta)<0$.

The map $\theta\mapsto h_\theta$ is continuous from the compact set $\Theta_{\theta_0}$ into
$C^2$ with the uniform norms, by dominated convergence in the Fourier integrals with the weights
$1$, $|t|$ and $t^2$, and $c_\theta$ is bounded below there. Since $x_\theta$ is the unique zero of
$h'_\theta$ and $h''_\theta(x_\theta)<0$, the implicit function theorem makes $x_\theta$
continuous. The last two claims follow from the continuity of $(\theta,x)\mapsto h''_\theta(x)$
and of $|h'_\theta(x)|$, which is positive on the compact set of the statement.
\end{proof}

Yamazato's unimodality of stable laws does not suffice here, since it gives neither
$h'_\theta\ne0$ away from the mode nor $h''_\theta<0$ at the mode.

\begin{theorem}[a local limit theorem for differences in the bulk]\label{thm:bulk-diff}
Assume (SPAN), (LR) and (LD2) for both signs and the conditions of Section~\ref{app:zone}, and
let the start be fresh with either sign or fair. For every $R$, uniformly in $|x_{N,k}|\le R$,
\[
 B_N^2\Delta P_N(k)=h'_N(x_{N,k})+o(1),\qquad B_N^3\Delta^2P_N(k)=h''_N(x_{N,k})+o(1).
\]
\end{theorem}

\begin{proof}
By linearity we take a fresh start of sign $\sgn_0$, so $P_N=F^{\sgn_0}_N$. We follow Step~5 of
the proof of Theorem~\ref{thm:bulk}(a) and keep its notation. For large $N$ the points with
$|x_{N,k}|\le R$ satisfy $\varepsilon N\le k\le(1-\varepsilon)N$. For one product
$\phi(k)=A(k)B(N-k)$ in~\eqref{eq:bulk-pairs}, with $A$ the plus kernel and $B$ the minus kernel,
\begin{align*}
 \Delta\phi(k)&=\Delta A(k)\,B(N-k-1)-A(k)\,\Delta B(N-k-1),\\
 \Delta^2\phi(k)&=\Delta^2A(k)\,B(N-k-2)-2\Delta A(k)\,\Delta B(N-k-2)+A(k)\,\Delta^2B(N-k-2).
\end{align*}
So a difference on the factor of sign $\sgn$ carries the sign $\varsigma_\sgn$. Fix $K\ge2$ and let
$I=I^{\bar\tau}_{k_{\bar\tau}}(K)$.

Indices in $I$. Here $a_\tau(m)=B_N(1+o(1))$ and $a_{\bar\tau}(m)=b_{\bar\tau,N}(1+o(1))$, with
$b_{\bar\tau,N}=\theta_{\bar\tau,N}B_N\ge\theta_0B_N$. By Lemma~\ref{lem:fourier-diff} and the
uniform continuity of $g^{(r)}_\sgn$, uniformly in $m\in I$ and in shifts of the levels by at most
$2$,
\[
 \Delta^{r_1}\mathsf k^\tau_m(k_\tau)=\frac{c\,g^{(r_1)}_\tau(\zeta_{N,x}(y_m))+o(1)}{B_N^{1+r_1}},\qquad
 \Delta^{r_2}\mathsf k^{\bar\tau}_m(k_{\bar\tau})=\frac{c'\,g^{(r_2)}_{\bar\tau}(y_m)+o(1)}{b_{\bar\tau,N}^{1+r_2}},
\]
with $c,c'$ the weights of the two kernels. The points $y_m$ form a progression of step
$\mu_{\bar\tau}/b_{\bar\tau,N}\cdot(1+o(1))$ in $[-K,K]$, with at most $Cb_{\bar\tau,N}$ points. So the
sum over $m\in I$ of the product is a Riemann sum,
\[
 \sum_{m\in I}\Delta^{r_1}\mathsf k^\tau_m\,\Delta^{r_2}\mathsf k^{\bar\tau}_m
 =\frac{cc'}{\mu_{\bar\tau}B_N^{1+r_1}b_{\bar\tau,N}^{r_2}}\int_{-K}^K g^{(r_1)}_\tau(\zeta_{N,x}(y))\,g^{(r_2)}_{\bar\tau}(y)\,dy
 +o(B_N^{-1-r_1-r_2}),
\]
where we used $b_{\bar\tau,N}\ge\theta_0B_N$ for the error. Extending the integral to $\R$ costs at
most $C\int_{|y|>K}|g^{(r_2)}_{\bar\tau}|$. Put $\Phi(z,y)=g_\tau(z)g_{\bar\tau}(y)$ and
$\theta=\theta_{\bar\tau,N}$. Summing over the products and over $i$ with~\eqref{eq:bulk-weights},
the main term of $B_N^{1+r}\Delta^rP_N(k)$ is
\[
 \frac s{\mu_{\bar\tau}}\int_\R\Bigl[\Bigl(\varsigma_\tau\partial_1+\frac{\varsigma_{\bar\tau}}\theta\partial_2\Bigr)^r\Phi\Bigr](\zeta_{N,x}(y),y)\,dy .
\]
Since $\zeta_{N,x}$ has the slope $\zeta'=\mu_\tau\theta/\mu_{\bar\tau}$, every integrable $\Psi$
with integrable derivatives satisfies $\int\partial_2\Psi(\zeta_{N,x}(y),y)dy=-\zeta'\int\partial_1\Psi(\zeta_{N,x}(y),y)dy$.
Applied once or twice, this turns the operator into
$(\varsigma_\tau-\varsigma_{\bar\tau}\zeta'/\theta)^r\partial_1^r=(\varsigma_\tau s/\mu_{\bar\tau})^r\partial_1^r$,
since $\varsigma_{\bar\tau}=-\varsigma_\tau$. By~\eqref{eq:bulk-profile} the result is
$h^{(r)}_N(x)$.

Indices outside $I$. If the differences fall only on the $\tau$ factor, then for
$m\ge m_\varepsilon$ Lemma~\ref{lem:tails}(b) and Lemma~\ref{lem:index}(a) give
$O(\eta(K))B_N^{-1-r}$, and for $m<m_\varepsilon$ Lemma~\ref{lem:tails}(c) gives $o(B_N^{-1-r})$.
If they fall only on the $\bar\tau$ factor, then $\mathsf k^\tau_m\le C/B_N$
by~\eqref{eq:bulk-factor}, and Lemma~\ref{lem:tails}(a) at the level $k_{\bar\tau}$, with
$a_{\bar\tau}(k_{\bar\tau}/\mu_{\bar\tau})\ge c\theta_0B_N$, gives $O(K^{1-\alpha+\delta}+e^{-cK})B_N^{-1-r}$.
In the mixed term for $r=2$ we combine Lemma~\ref{lem:tails}(b) for the $\tau$ factor with
Lemma~\ref{lem:tails}(a) for $r=1$ when $m\ge m_\varepsilon$, and use Lemma~\ref{lem:tails}(c)
when $m<m_\varepsilon$. Every product in~\eqref{eq:bulk-pairs} pairs a $u$ kernel with a $q$
kernel. Letting $N\to\infty$ and then $K\to\infty$ proves the theorem.
\end{proof}

\begin{corollary}[one change of sign in the bulk]\label{cor:bulk-sign}
Under the hypotheses of Theorem~\ref{thm:bulk-diff}, let $R_0=\sup_{\Theta_{\theta_0}}|x_\theta|$
and $R\ge R_0+1$. For all large $N$ the sign of $\Delta P_N(k)$ on $|x_{N,k}|\le R$ changes exactly
once, from positive to negative, at a point $k_N$ with $x_{N,k_N}-x_{\theta_N}\to0$, and $\Delta P_N$
is strictly decreasing on $|x_{N,k}-x_{\theta_N}|\le\varepsilon_0/2$.
\end{corollary}

\begin{proof}
Let $0<\varepsilon\le\varepsilon_0/2$. By Theorem~\ref{thm:bulk-diff} with $R+1$ and
Lemma~\ref{lem:bell}, for large $N$, $\Delta P_N(k)>0$ when $x_{N,k}\le x_{\theta_N}-\varepsilon$ and
$\Delta P_N(k)<0$ when $x_{N,k}\ge x_{\theta_N}+\varepsilon$, and $\Delta^2P_N(k)<0$ when
$|x_{N,k}-x_{\theta_N}|\le\varepsilon_0/2$, since $x_{N,k+2}-x_{N,k}=2/B_N\to0$. So on the window
around $x_{\theta_N}$ the sequence $\Delta P_N$ decreases strictly from a positive to a negative
value, and its sign changes exactly once, inside $|x-x_{\theta_N}|<\varepsilon$. Since
$\varepsilon$ is arbitrary, $x_{N,k_N}-x_{\theta_N}\to0$.
\end{proof}

\subsection{Junctions and the proof of Theorem~\ref{thm:mode-count}}\label{app:junctions}

\emph{(EQ) and comparable scales.} Under (EQ) the ratio $T_\plus/T_\minus$ is regularly varying
with index $\alpha_\minus-\alpha_\plus$ and bounded above and below, so
$\alpha_\plus=\alpha_\minus$. If $a_\minus(t)\ge Ca_\plus(t)$ with $C$ large, then by
Lemma~\ref{lem:rv}(ii)
$1\sim tT_\minus(a_\minus(t))\le tT_\minus(Ca_\plus(t))\le2C^{-\alpha/2}tT_\minus(a_\plus(t))\le c'C^{-\alpha/2}$,
which is impossible. So $a_\plus\asymp a_\minus$, and $\min_\sgn\theta_{\sgn,N}\ge\theta_0$ for all
large $N$ and some $\theta_0>0$. Hence Theorem~\ref{thm:mode-count} is the case of
Corollary~\ref{cor:sequences} in which every large $N$ qualifies, and it suffices to prove the
corollary. Its hypotheses are those under which Sections~\ref{app:zone} and~\ref{app:bulk} were
proved.

\emph{The zones.} Parts (i) and (ii) are Theorem~\ref{thm:ends}(a) and~(b), which need (RV),
(SPAN), (LR) and (LD) for both signs. (RV) follows from (LR) and (LD) from Lemma~\ref{lem:ld2}.
Part (iii) was proved in Section~\ref{app:antimode} and part (iv) in Section~\ref{app:zone}. For
(v) take $K_0$ as in the proof of (iv), enlarged if necessary so that $K_0\ge R_0$, and apply Corollary~\ref{cor:bulk-sign} with
$R=K_0+1$. The mirror statements at the minus end follow from the same proofs with the signs
exchanged, using (LD2) for the minus sign.

\emph{The junctions.} The ranges overlap as follows.
\begin{itemize}[nosep]
\item Parts (i) and (ii) meet at $J_\plus$.
\item Part (ii) ends at $(x_*-\eta')\Lambda_N$ and part (iv) starts at $(x_*+\eta')\Lambda_N$, and the
window of (iii) is the open interval between them. By the proof of Theorem~\ref{thm:ends}(d),
$D_N<0$ at the first integer of the window and $D_N>0$ at its last integer.
\item Part (iv) is the union of Theorem~\ref{thm:ends}(c) up to $N^{1-\eta}$ and
Theorem~\ref{thm:zone-diff} from $N^{1-\eta}$ on, with the same $\eta<1-1/\alpha$.
\item Part (iv) covers the points with $x_{N,N-j}\ge K_0$, since $j\le(1-p)N-K_0B_N$ is
$N-j\ge pN+K_0B_N$. Part (v) covers $|x_{N,k}|\le K_0+1$. So the two overlap.
\item The minus end is the mirror image, and covers $x_{N,k}\le-K_0$.
\end{itemize}
The constants are chosen in the order $\eta,\eta'$, then $K_0$, which therefore depends on the
laws and on $\eta$ and $\eta'$, then $N$. Each zone holds for
all large $N$, so all hold together for all large $N$, and the proof of the count in
Section~\ref{sec:mode-count} applies. This proves Corollary~\ref{cor:sequences} and
Theorem~\ref{thm:mode-count}. \qed
